\documentclass[12pt]{article}
\usepackage[T1]{fontenc}
\usepackage[utf8]{inputenc}
\usepackage{lmodern}
\usepackage{textcomp}
\usepackage{enumitem}
\usepackage{geometry}
\geometry{margin=1in}
\begin{document}
\begin{center}
Answer Key - Sociology - Undergraduate Level Assessment 21 \
\small (Answers are ordered exactly as in the question paper.)
\end{center}

\section*{Multiple Choice}
\begin{enumerate}[label=\arabic*.]
\item \textbf{Answer:} B
\\ \textit{Options:} A) Westernization theory; B) Modernization theory; C) Industrialization theory; D) Dependency theory
\\ \textit{Note:} Correct option: B - Modernization theory
\item \textbf{Answer:} A
\\ \textit{Options:} A) those in non-manual occupations, exercising authority on behalf of the state; B) people working in consultancy firms who were recruited by big businesses; C) the young men and women employed in domestic service in the nineteenth century; D) those who had worked in the armed services
\\ \textit{Note:} Correct option: A - those in non-manual occupations, exercising authority on behalf of the state
\item \textbf{Answer:} A
\\ \textit{Options:} A) we internalise and take on social roles from a pre-existing framework; B) we create and negotiate our roles through interaction with others; C) social roles are not fixed or stable but fluid and pluralistic; D) roles have to be learned to suppress unconscious motivations
\\ \textit{Note:} Correct option: A - we internalise and take on social roles from a pre-existing framework
\item \textbf{Answer:} C
\\ \textit{Options:} A) the findings are amenable to statistical analysis; B) it is conducted over a period of several years; C) it uncovers rich, detailed accounts from an insider's perspective; D) it compares findings from a number of different cases
\\ \textit{Note:} Correct option: C - it uncovers rich, detailed accounts from an insider's perspective
\item \textbf{Answer:} C
\\ \textit{Options:} A) genetics; B) evolution; C) height; D) brain size
\\ \textit{Note:} Correct option: C - height
\end{enumerate}

\section*{True / False}
\begin{enumerate}[label=\arabic*.]
\setcounter{enumi}{5}
\item \textbf{Answer:} True
\\ \textit{Options:} A) True; B) False
\\ \textit{Note:} LLM-generated true statement based on MCQ.
\item \textbf{Answer:} False
\\ \textit{Options:} A) True; B) False
\\ \textit{Note:} LLM-generated false statement. Correct answer: 'all of the above'.
\item \textbf{Answer:} True
\\ \textit{Options:} A) True; B) False
\\ \textit{Note:} LLM-generated true statement based on MCQ.
\item \textbf{Answer:} False
\\ \textit{Options:} A) True; B) False
\\ \textit{Note:} LLM-generated false statement. Correct answer: 'a minimum wage'.
\item \textbf{Answer:} True
\\ \textit{Options:} A) True; B) False
\\ \textit{Note:} LLM-generated true statement based on MCQ.
\end{enumerate}

\section*{Long Form Response}
\begin{enumerate}[label=\arabic*.]
\setcounter{enumi}{10}
\item \textbf{Answer:} Wider at the bottom than at the top. Explain why this is the correct answer and provide supporting evidence. Reference key facts or concepts that support this choice.
\\ \textit{Note:} Long-form derived from MCQ; award credit for stating the correct idea and giving supporting reasoning.
\item \textbf{Answer:} most strikes go unnoticed by employers and the mass media. Explain why this is the correct answer and provide supporting evidence. Reference key facts or concepts that support this choice.
\\ \textit{Note:} Long-form derived from MCQ; award credit for stating the correct idea and giving supporting reasoning.
\end{enumerate}

\vfill
\noindent\textit{End of Answer Key}
\end{document}